\chapter{2016年全国III卷（理）}

\section{单选题}

\begin{problem}
设集合 \(S=\{x\mid (x-2)(x-3)\ge0\}\)，\(T=\{x\mid x>0\}\)，则 \(S\cap T=\)（\quad）

\choices
    {\([2,3]\)}
    {\((-\infty,2]\cup[3,+\infty)\)}
    {\([3,+\infty)\)}
    {\((0,2]\cup[3,+\infty)\)}
\end{problem}

\begin{answer}
D
\end{answer}

\begin{solution}
由
\[
(x-2)(x-3)\ge0
\]
得 \(x\le2\) 或 \(x\ge3\)，所以
\[
S=(-\infty,2]\cup[3,+\infty)
\]
再与 \(T=(0,+\infty)\) 取交集，得到
\[
S\cap T=(0,2]\cup[3,+\infty)
\]
因此选择 D.
\end{solution}

\begin{problem}
若 \(z=1+2\i\)，则 \(\dfrac{4\i}{z\cdot \overline z-1}=\)（\quad）

\choices
    {\(1\)}
    {\(-1\)}
    {\(\i\)}
    {\(-\i\)}
\end{problem}

\begin{answer}
C
\end{answer}

\begin{solution}
因为 \(\overline z=1-2\i\)，所以
\[
z\overline z=(1+2\i)(1-2\i)=1+4=5
\]
于是
\[
\frac{4\i}{z\overline z-1}=\frac{4\i}{5-1}=\i
\]
因此选择 C.
\end{solution}

\begin{problem}
已知向量 \(\overrightarrow{BA}=\left(\dfrac12,\dfrac{\sqrt3}{2}\right)\)，\(\overrightarrow{BC}=\left(\dfrac{\sqrt3}{2},\dfrac12\right)\)，则 \(\angle ABC=\)（\quad）

\choices
    {\(30^\circ\)}
    {\(45^\circ\)}
    {\(60^\circ\)}
    {\(120^\circ\)}
\end{problem}

\begin{answer}
A
\end{answer}

\begin{solution}
两向量的长度均为
\[
|\overrightarrow{BA}|=|\overrightarrow{BC}|=\sqrt{\frac14+\frac34}=1
\]
它们的数量积为
\[
\overrightarrow{BA}\cdot\overrightarrow{BC}
=\frac12\cdot\frac{\sqrt3}{2}+\frac{\sqrt3}{2}\cdot\frac12
=\frac{\sqrt3}{2}
\]
因此
\[
\cos\angle ABC
=\frac{\overrightarrow{BA}\cdot\overrightarrow{BC}}
{|\overrightarrow{BA}|\,|\overrightarrow{BC}|}
=\frac{\sqrt3}{2}
\]
由于 \(0\le\angle ABC\le\pi\)，得 \(\angle ABC=30^\circ\)，所以选择 A.
\end{solution}

\begin{problem}
某旅游城市为向游客介绍本地的气温情况，绘制了一年中月平均最高气温和平均最低气温的雷达图. 图中 \(A\) 点表示十月的平均最高气温约为 \(15^\circ\mathrm{C}\)，\(B\) 点表示四月的平均最低气温约为 \(5^\circ\mathrm{C}\). 下面叙述不正确的是（\quad）

\begingroup
\bitmapfigure[width=0.92\linewidth]{img/2016/national_paper_3_science/q4_fig1.png}
\endgroup

\choices
    {各月的平均最低气温都在 \(0^\circ\mathrm{C}\) 以上}
    {七月的平均温差比一月的平均温差大}
    {三月和十一月的平均最高气温基本相同}
    {平均气温高于 \(20^\circ\mathrm{C}\) 的月份有 \(5\) 个}
\end{problem}

\begin{answer}
D
\end{answer}

\begin{solution}
逐月比较雷达图中的虚线和实线可知，各月平均最低气温均在 \(0^\circ\mathrm{C}\) 以上；七月两条曲线的径向差距大于一月；三月和十一月的平均最高气温均约为 \(10^\circ\mathrm{C}\)，基本相同.

再看平均最高气温的实线，超过 \(20^\circ\mathrm{C}\) 的月份只有七月和八月，六月约为 \(20^\circ\mathrm{C}\)，并非有 \(5\) 个. 因此不正确的叙述是 D.
\end{solution}

\begin{problem}
若 \(\tan\alpha=\dfrac34\)，则 \(\cos^2\alpha+2\sin 2\alpha=\)（\quad）

\choices
    {\(\dfrac{64}{25}\)}
    {\(\dfrac{48}{25}\)}
    {\(1\)}
    {\(\dfrac{16}{25}\)}
\end{problem}

\begin{answer}
A
\end{answer}

\begin{solution}
由 \(1+\tan^2\alpha=\sec^2\alpha\)，得
\[
\cos^2\alpha=\frac{1}{1+\tan^2\alpha}
=\frac{1}{1+\frac{9}{16}}=\frac{16}{25}
\]
又
\[
\sin2\alpha=\frac{2\tan\alpha}{1+\tan^2\alpha}
=\frac{2\cdot\frac34}{1+\frac{9}{16}}=\frac{24}{25}
\]
所以
\[
\cos^2\alpha+2\sin2\alpha
=\frac{16}{25}+2\cdot\frac{24}{25}
=\frac{64}{25}
\]
因此选择 A.
\end{solution}

\begin{problem}
已知 \(a=2^{4/3}\)，\(b=4^{2/5}\)，\(c=25^{1/3}\)，则（\quad）

\choices
    {\(b<a<c\)}
    {\(a<b<c\)}
    {\(b<c<a\)}
    {\(c<a<b\)}
\end{problem}

\begin{answer}
A
\end{answer}

\begin{solution}
先把 \(b\) 化为以 \(2\) 为底的幂：
\[
b=4^{2/5}=2^{4/5}
\]
由于 \(2>1\)，且 \(\frac45<\frac43\)，所以 \(b<a\).

再比较 \(a\) 与 \(c\)：
\[
a^3=2^4=16<25=c^3
\]
由于 \(a\)，\(c>0\)，所以 \(a<c\). 因而 \(b<a<c\)，选择 A.
\end{solution}

\begin{problem}
执行下图的程序框图，如果输入的 \(a=4\)，\(b=6\)，那么输出的 \(n=\)（\quad）

\begingroup
\bitmapfigure[width=0.42\linewidth]{img/2016/national_paper_3_science/q7_fig1.png}
\endgroup

\choices
    {\(3\)}
    {\(4\)}
    {\(5\)}
    {\(6\)}
\end{problem}

\begin{answer}
B
\end{answer}

\begin{solution}
按照框图中三个赋值语句的先后顺序逐次循环. 初始 \(a=4\)，\(b=6\)，\(n=0\)，\(s=0\).

第一次循环后，
\(a=6-4=2\)，\(b=6-2=4\)，\(a=4+2=6\)
从而 \(s=6\)，\(n=1\).

第二次循环后，
\(a=4-6=-2\)，\(b=4-(-2)=6\)，\(a=6+(-2)=4\)
从而 \(s=10\)，\(n=2\).

第三次循环后，
\(a=6-4=2\)，\(b=6-2=4\)，\(a=4+2=6\)
从而 \(s=16\)，\(n=3\). 此时 \(s>16\) 不成立.

第四次循环后，
\(a=4-6=-2\)，\(b=4-(-2)=6\)，\(a=6+(-2)=4\)
从而 \(s=20\)，\(n=4\). 此时 \(s>16\) 成立，输出 \(n=4\)，所以选择 B.
\end{solution}

\begin{problem}
在 \(\triangle ABC\) 中，\(B=\dfrac{\pi}{4}\)，\(BC\) 边上的高等于 \(\dfrac13BC\)，则 \(\cos A=\)（\quad）

\choices
    {\(\dfrac{3\sqrt{10}}{10}\)}
    {\(\dfrac{\sqrt{10}}{10}\)}
    {\(-\dfrac{\sqrt{10}}{10}\)}
    {\(-\dfrac{3\sqrt{10}}{10}\)}
\end{problem}

\begin{answer}
C
\end{answer}

\begin{solution}
设从 \(A\) 向 \(BC\) 作高，垂足为 \(D\)，并令 \(AD=h>0\). 由题意 \(BC=3h\). 因为 \(\angle B=\frac\pi4\)，在直角三角形 \(ABD\) 中
\[
BD=AD\cot\frac\pi4=h
\]
所以 \(D\) 在线段 \(BC\) 上，且 \(DC=BC-BD=2h\). 于是
\(AB=\sqrt{AD^2+BD^2}=h\sqrt2\)，\(AC=\sqrt{AD^2+DC^2}=h\sqrt5\).
由余弦定理，
\[
\cos A=\frac{AB^2+AC^2-BC^2}{2AB\cdot AC}
=\frac{2h^2+5h^2-9h^2}{2h\sqrt2\cdot h\sqrt5}
=-\frac{1}{\sqrt{10}}=-\frac{\sqrt{10}}{10}
\]
因此选择 C.
\end{solution}

\begin{problem}
如图，网格纸上小正方形的边长为 \(1\)，粗实线画出的是某多面体的三视图，则该多面体的表面积为（\quad）

\begingroup
\bitmapfigure[width=0.68\linewidth]{img/2016/national_paper_3_science/q9_fig1.png}
\endgroup

\choices
    {\(18+36\sqrt5\)}
    {\(54+18\sqrt5\)}
    {\(90\)}
    {\(81\)}
\end{problem}

\begin{answer}
B
\end{answer}

\begin{solution}
由三视图还原可知，该多面体是斜四棱柱，可取图中的平行四边形为底面. 由网格图读出其一条底边长为 \(3\)，对应的高为 \(6\)，所以底面面积为

\[
S_\text{底}=3\times6=18
\]

该平行四边形的另一条边长为

\[
\sqrt{3^2+6^2}=3\sqrt5
\]
，另外两个投影显示两底面之间的棱长为 \(3\). 因此四个侧面中有两个面积为 \(3\times3\) 的矩形，另两个面积为 \(3\times3\sqrt5\) 的平行四边形.

因此表面积为
\[
2S_\text{底}+2\left(3\times3+3\times3\sqrt5\right)
=36+18+18\sqrt5=54+18\sqrt5
\]
所以选择 B.
\end{solution}

\begin{problem}
在封闭的直三棱柱 \(ABC-A_1B_1C_1\) 内有一个体积为 \(V\) 的球，若 \(AB\perp BC\)，\(AB=6\)，\(BC=8\)，\(AA_1=3\)，则 \(V\) 的最大值是（\quad）

\choices
    {\(4\pi\)}
    {\(\dfrac{9\pi}{2}\)}
    {\(6\pi\)}
    {\(\dfrac{32\pi}{3}\)}
\end{problem}

\begin{answer}
B
\end{answer}

\begin{solution}
底面三角形 \(ABC\) 是直角三角形，且
\[
AC=\sqrt{AB^2+BC^2}=\sqrt{6^2+8^2}=10
\]
其内切圆半径为
\[
r=\frac{AB+BC-AC}{2}=\frac{6+8-10}{2}=2
\]
球的半径既不能超过直三棱柱高的一半，也不能超过底面三角形内切圆半径，因此
\[
R\le\min\left(\frac{AA_1}{2},r\right)
=\min\left(\frac32,2\right)=\frac32
\]
取球心在棱柱高度的中点、在底面上的投影为三角形内心时可达到 \(R=\frac32\)，所以
\[
V_{\max}=\frac43\pi\left(\frac32\right)^3=\frac{9\pi}{2}
\]
因此选择 B.
\end{solution}

\begin{problem}
已知 \(O\) 为坐标原点，\(F\) 是椭圆 \(C:\dfrac{x^2}{a^2}+\dfrac{y^2}{b^2}=1\)（\(a>b>0\)）的左焦点，\(A\)，\(B\) 分别为 \(C\) 的左，右顶点，\(P\) 为 \(C\) 上一点，且 \(PF\perp x\) 轴，过点 \(A\) 的直线 \(l\) 与线段 \(PF\) 交于点 \(M\)，与 \(y\) 轴交于点 \(E\). 若直线 \(BM\) 经过 \(OE\) 的中点，则 \(C\) 的离心率为（\quad）

\choices
    {\(\dfrac13\)}
    {\(\dfrac12\)}
    {\(\dfrac23\)}
    {\(\dfrac34\)}
\end{problem}

\begin{answer}
A
\end{answer}

\begin{solution}
设椭圆焦距的一半为 \(c=\sqrt{a^2-b^2}\)，则
\(F=(-c,0)\)，\(A=(-a,0)\)，\(B=(a,0)\).
由于 \(PF\perp x\) 轴，点 \(P\) 的横坐标为 \(-c\). 取 \(P\) 在 \(x\) 轴上方的情形，另一情形关于 \(x\) 轴对称. 按通常的非退化构型取 \(M\ne F\)，设直线 \(l\) 的斜率为 \(k\)，则 \(k\ne0\)，且
\(M=(-c,k(a-c))\)，\(E=(0,ka)\).
令 \(N\) 为 \(OE\) 的中点，则 \(N=(0,ka/2)\). 由 \(B\)，\(M\)，\(N\) 三点共线，比较直线 \(BN\) 与 \(BM\) 的斜率，得
\[
-\frac{k}{2}=-\frac{k(a-c)}{a+c}
\]
约去 \(-k\)，得到
\(\frac12=\frac{a-c}{a+c}\)，\(a+c=2a-2c\)，\(a=3c\).
所以椭圆的离心率为
\[
e=\frac ca=\frac13
\]
因此选择 A.
\end{solution}

\begin{problem}
定义“规范 \(01\) 数列”\(\{a_n\}\) 如下：\(\{a_n\}\) 共有 \(2m\) 项，其中 \(m\) 项为 \(0\)，\(m\) 项为 \(1\)，且对任意 \(k\le2m\)，\(a_1\)，\(a_2\)，\dots，\(a_k\) 中 \(0\) 的个数不少于 \(1\) 的个数. 若 \(m=4\)，则不同的“规范 \(01\) 数列”共有（\quad）

\choices
    {\(18\) 个}
    {\(16\) 个}
    {\(14\) 个}
    {\(12\) 个}
\end{problem}

\begin{answer}
C
\end{answer}

\begin{solution}
不计限制时，长度为 \(8\) 且含四个 \(0\)、四个 \(1\) 的数列共有
\[
\binom84=70
\]
个.

考虑不规范的数列. 对每个不规范数列，取第一次出现“前缀中 \(1\) 的个数比 \(0\) 的个数多 \(1\)”的位置，并将此前缀中的 \(0\) 与 \(1\) 互换. 这样得到一个含有五个 \(1\)、三个 \(0\) 的长度为 \(8\) 数列；反过来，对任意含五个 \(1\)、三个 \(0\) 的数列，在第一次使 \(1\) 的个数比 \(0\) 的个数多 \(1\) 之前作同样互换，可唯一还原一个不规范数列.

因此不规范数列共有
\[
\binom83=56
\]
个，规范数列共有
\[
70-56=14
\]
个. 所以选择 C.
\end{solution}

\section{填空题}

\begin{problem}
若 \(x\)，\(y\) 满足约束条件
\(\begin{cases}
x-y+1\ge0,\\
x-2y\le0,\\
x+2y-2\le0,
\end{cases}\)
则 \(z=x+y\) 的最大值为\fillinblank{}.
\end{problem}

\begin{answer}
\(\dfrac{3}{2}\)
\end{answer}

\begin{solution}
三个约束分别给出
\(y\le x+1\)，\(y\ge\frac{x}{2}\)，\(y\le1-\frac{x}{2}\).
可行域是三条边界直线围成的三角形，其顶点由两两联立得到：
\(y=x+1\)，\(\ y=\frac{x}{2}\Longrightarrow(-2,-1)\)，\(y=x+1\)，\(\ y=1-\frac{x}{2}\Longrightarrow(0,1)\)，\(y=\frac{x}{2}\)，\(\ y=1-\frac{x}{2}\Longrightarrow(1,\tfrac12)\).
目标函数 \(z=x+y\) 在线性可行域上的最大值在顶点处取得. 三个顶点对应的 \(z\) 值分别为 \(-3\)，\(1\)，\(\frac32\)，所以
\[
z_{\max}=\frac32
\]
\end{solution}

\begin{problem}
函数 \(y=\sin x-\sqrt3\cos x\) 的图象可由函数 \(y=\sin x+\sqrt3\cos x\) 的图象至少向右平移\fillinblank{} 个单位长度得到.
\end{problem}

\begin{answer}
\(\dfrac{2\pi}{3}\)
\end{answer}

\begin{solution}
利用辅助角公式，
\[
\sin x+\sqrt3\cos x=2\sin\left(x+\frac\pi3\right)
\]
\[
\sin x-\sqrt3\cos x=2\sin\left(x-\frac\pi3\right)
\]
若将前一个函数的图象向右平移 \(h\) 个单位，所得函数为
\[
2\sin\left(x-h+\frac\pi3\right)
\]
要与 \(2\sin(x-\frac\pi3)\) 相同，需要
\(-h+\frac\pi3=-\frac\pi3+2k\pi\)，\(k\in\mathbb Z\)
即 \(h=\frac{2\pi}{3}-2k\pi\). 最小的正平移量为
\[
h=\frac{2\pi}{3}
\]
\end{solution}

\begin{problem}
已知 \(f(x)\) 为偶函数，当 \(x<0\) 时，\(f(x)=\ln(-x)+3x\)，则曲线 \(y=f(x)\) 在点 \((1,-3)\) 处的切线方程是\fillinblank{}.
\end{problem}

\begin{answer}
\(2x+y+1=0\)
\end{answer}

\begin{solution}
因为 \(f\) 是偶函数，所以当 \(x>0\) 时
\[
f(x)=f(-x)=\ln x-3x
\]
于是 \(f(1)=-3\)，且
\(f'(x)=\frac1x-3\)，\(f'(1)=-2\).
曲线在 \((1,-3)\) 处的切线为
\[
y+3=-2(x-1)
\]
整理得
\[
2x+y+1=0
\]
\end{solution}

\begin{problem}
已知直线 \(l:mx+y+3m-\sqrt3=0\) 与圆 \(x^2+y^2=12\) 交于 \(A\)，\(B\) 两点，过 \(A\)，\(B\) 分别做 \(l\) 的垂线与 \(x\) 轴交于 \(C\)，\(D\) 两点，若 \(AB=2\sqrt3\)，则 \(|CD|=\fillinblank{}\).
\end{problem}

\begin{answer}
\(4\)
\end{answer}

\begin{solution}
由弦长公式 \(2\sqrt3=2\sqrt{12-d^2}\)，得圆心到 \(l\) 的距离 \(d=3\). 又
\(d=\frac{|3m-\sqrt3|}{\sqrt{m^2+1}}\)，\((3m-\sqrt3)^2=9(m^2+1)\)
所以 \(9m^2-6\sqrt3m+3=9m^2+9\)，即 \(m=-\frac1{\sqrt3}\). 直线 \(l\) 的斜率为 \(\frac1{\sqrt3}\)，因而弦 \(AB\) 的单位方向向量可取 \(\left(\frac{\sqrt3}{2},\frac12\right)\). 结合 \(AB=2\sqrt3\)，得
\[
\overrightarrow{AB}=\pm(3,\sqrt3)
\]
过圆上任意点 \((x,y)\) 作 \(l\) 的垂线，其斜率为 \(-\sqrt3\)，令其与 \(x\) 轴交于 \((u,0)\)，则 \(0-y=-\sqrt3(u-x)\)，从而 \(u=x+\frac{y}{\sqrt3}\). 因此两交点对应的横坐标之差的绝对值为
\[
|CD|=\left|3+\frac{\sqrt3}{\sqrt3}\right|=4
\]
\end{solution}

\section{解答题}

\begin{problem}
已知数列 \(\{a_n\}\) 的前 \(n\) 项和 \(S_n=1+\lambda a_n\)，其中 \(\lambda\ne0\)；

\begin{enumerate}
\item 证明 \(\{a_n\}\) 是等比数列，并求其通项公式；
\item 若 \(S_5=\dfrac{31}{32}\)，求 \(\lambda\).
\end{enumerate}
\end{problem}

\begin{solution}
\begin{enumerate}
\item 由 \(S_1=a_1=1+\lambda a_1\)，得 \((1-\lambda)a_1=1\)，所以 \(\lambda\ne1\)，且 \(a_1=\dfrac{1}{1-\lambda}\). 当 \(n\ge2\) 时，\(S_n=S_{n-1}+a_n\)，从而
\[
1+\lambda a_n=1+\lambda a_{n-1}+a_n
\]
因此 \((\lambda-1)a_n=\lambda a_{n-1}\)，即
\[
a_n=\frac{\lambda}{\lambda-1}a_{n-1}
\]
所以 \(\{a_n\}\) 是首项为 \(\dfrac{1}{1-\lambda}\)、公比为 \(\dfrac{\lambda}{\lambda-1}\) 的等比数列，通项为
\[
a_n=\frac{1}{1-\lambda}\left(\frac{\lambda}{\lambda-1}\right)^{n-1}
\]

\item 令 \(q=\dfrac{\lambda}{\lambda-1}\)，由上式可得 \(\lambda a_5=-q^5\)，故
\[
S_5=1+\lambda a_5=1-q^5
\]
由 \(S_5=\dfrac{31}{32}\)，得 \(q^5=\dfrac{1}{32}\)，所以 \(q=\dfrac12\). 因而 \(\dfrac{\lambda}{\lambda-1}=\dfrac12\)，解得 \(\lambda=-1\).
\end{enumerate}
\end{solution}

\begin{problem}
下图是我国 \(2008\) 年至 \(2014\) 年生活垃圾无害化处理量（单位：亿吨）的折线图.

注：年份代码 \(1\) 到 \(7\) 分别对应年份 \(2008\) 到 \(2014\).

参考数据：\(\sum_{i=1}^{7}y_i=9.32\)，\(\sum_{i=1}^{7}t_iy_i=40.17\)，\(\sqrt{\sum_{i=1}^{7}\left(y_i-\bar y\right)^2}=0.55\)，\(\sqrt7\approx2.646\).

参考公式：相关系数 \(r=\dfrac{\sum_{i=1}^{n}\left(t_i-\bar t\right)\left(y_i-\bar y\right)}{\sqrt{\sum_{i=1}^{n}\left(t_i-\bar t\right)^2\sum_{i=1}^{n}\left(y_i-\bar y\right)^2}}\)，回归方程 \(\hat y=\hat a+\hat b t\) 中斜率和截距的最小二乘估计公式分别为：\(\hat b=\dfrac{\sum_{i=1}^{n}\left(t_i-\bar t\right)\left(y_i-\bar y\right)}{\sum_{i=1}^{n}\left(t_i-\bar t\right)^2}\)，\(\hat a=\bar y-\hat b\bar t\).

\begin{enumerate}
\item 由折线图看出，可用线性回归模型拟合 \(y\) 与 \(t\) 的关系，请用相关系数加以说明；
\item 建立 \(y\) 关于 \(t\) 的回归方程（系数精确到 \(0.01\)），预测 \(2016\) 年我国生活垃圾无害化处理量.
\end{enumerate}

\begingroup
\FigureLayoutDeclare{img/2016/national_paper_3_science/q18_fig1.png}{0.62\linewidth}{0bp 0bp 0bp 0bp}
\bitmapfigure[float=inline,max width=0.62\linewidth]{img/2016/national_paper_3_science/q18_fig1.png}
\endgroup
\end{problem}

\begin{solution}
\begin{enumerate}
\item 由 \(t_i=1\)，\(2\)，\(\ldots\)，\(7\)，得 \(\bar t=4\)，且
\[
\sum_{i=1}^{7}(t_i-\bar t)^2=9+4+1+0+1+4+9=28
\]
又因为 \(\sum_{i=1}^{7}(t_i-\bar t)=0\)，所以
\[
\sum_{i=1}^{7}(t_i-\bar t)(y_i-\bar y)=\sum_{i=1}^{7}t_i y_i-\bar t\sum_{i=1}^{7}y_i=40.17-4\times9.32=2.89
\]
因此
\[
r=\frac{2.89}{\sqrt{28}\times0.55}=\frac{2.89}{2\sqrt7\times0.55}\approx0.99
\]
相关系数接近 \(1\)，说明 \(y\) 与 \(t\) 正相关性很强，用线性回归模型拟合是合适的.

\item 由 \(\bar y=\dfrac{9.32}{7}\)，得
\(\hat b=\frac{2.89}{28}\approx0.10\)，\(\hat a=\bar y-\hat b\bar t=\frac{9.32}{7}-4\times\frac{2.89}{28}\approx0.92\).
所以回归方程为 \(\hat y=0.92+0.10t\). 年份 \(2016\) 对应的年份代码为 \(t=9\)，代入回归方程得
\[
\hat y=0.92+0.10\times9=1.82
\]
故预测 \(2016\) 年我国生活垃圾无害化处理量约为 \(1.82\) 亿吨.
\end{enumerate}
\end{solution}

\begin{problem}
如图，四棱锥 \(P-ABCD\) 中，\(PA\perp\) 底面 \(ABCD\)，\(AD\parallel BC\)，\(AB=AD=AC=3\)，\(PA=BC=4\)，\(M\) 为线段 \(AD\) 上一点，\(AM=2MD\)，\(N\) 为 \(PC\) 的中点.

\begin{enumerate}
\item 证明 \(MN\parallel\) 平面 \(PAB\)；
\item 求直线 \(AN\) 与平面 \(PMN\) 所成角的正弦值.
\end{enumerate}

\begingroup
\FigureLayoutDeclare{img/2016/national_paper_3_science/q19_fig1.png}{11.0cm}{531bp 97bp 221bp 48bp}
\bitmapfigure[float=inline,max width=\linewidth]{img/2016/national_paper_3_science/q19_fig1.png}
\endgroup
\end{problem}

\begin{solution}
在底面内建立直角坐标系，并取
\(A=(0,0,0)\)，\(D=(3,0,0)\)，\(B=(-2,\sqrt5,0)\)，\(C=(2,\sqrt5,0)\).
这是因为 \(AD\parallel BC\)、\(AD=3\)、\(BC=4\)，又由 \(AB=AC=3\) 可确定 \(B\)，\(C\) 的上述坐标. 由 \(PA\perp\) 底面且 \(PA=4\)，可取 \(P=(0,0,4)\). 另外，\(AM=2\)、\(N\) 为 \(PC\) 的中点，所以
\(M=(2,0,0)\)，\(N=\left(1,\frac{\sqrt5}{2},2\right)\).

\begin{enumerate}
\item 有
\[
\overrightarrow{MN}=\left(-1,\frac{\sqrt5}{2},2\right)=\frac12\overrightarrow{AB}+\frac12\overrightarrow{AP}
\]
因此 \(\overrightarrow{MN}\) 平行于平面 \(PAB\). 进一步，平面 \(PAB\) 的方程为 \(\sqrt5x+2y=0\)，而 \(M=(2,0,0)\) 不满足该方程，所以 \(M\notin\) 平面 \(PAB\). 因此直线 \(MN\) 不在且不交平面 \(PAB\)，从而 \(MN\parallel\) 平面 \(PAB\).

\item 取
\(\overrightarrow{PM}=(2,0,-4)\)，\(\overrightarrow{PN}=\left(1,\frac{\sqrt5}{2},-2\right)\).
它们的一个法向量为
\[
\overrightarrow{PM}\times\overrightarrow{PN}=(2\sqrt5,0,\sqrt5)
\]
故可取平面 \(PMN\) 的法向量为 \(\bs{n}=(2,0,1)\). 又
\(\overrightarrow{AN}=\left(1,\frac{\sqrt5}{2},2\right)\)，\(\left|\overrightarrow{AN}\right|=\frac52\)，\(|\bs{n}|=\sqrt5\).
设直线 \(AN\) 与平面 \(PMN\) 所成的角为 \(\theta\)，则
\[
\sin\theta=\frac{\left|\overrightarrow{AN}\cdot\bs{n}\right|}{\left|\overrightarrow{AN}\right||\bs{n}|}=\frac{|2+2|}{\frac52\sqrt5}=\frac{8\sqrt5}{25}
\]
\end{enumerate}
\end{solution}

\begin{problem}
已知抛物线 \(C:y^2=2x\) 的焦点为 \(F\)，平行于 \(x\) 轴的两条直线 \(l_1\)，\(l_2\) 分别交 \(C\) 于 \(A\)，\(B\) 两点，交 \(C\) 的准线于 \(P\)，\(Q\) 两点.

\begin{enumerate}
\item 若 \(F\) 在线段 \(AB\) 上，\(R\) 是 \(PQ\) 的中点，证明 \(AR\parallel FQ\)；
\item 若 \(\triangle PQF\) 的面积是 \(\triangle ABF\) 的面积的两倍，求 \(AB\) 中点的轨迹方程.
\end{enumerate}
\end{problem}

\begin{solution}
\begin{enumerate}
\item 设 \(l_1:y=a\)，\(l_2:y=b\)，则
\(A=\left(\frac{a^2}{2},a\right)\)，\(B=\left(\frac{b^2}{2},b\right)\)，\(P=\left(-\frac12,a\right)\)，\(Q=\left(-\frac12,b\right)\)，\(F=\left(\frac12,0\right)\).
直线 \(AB\) 的方程为
\[
2x-(a+b)y+ab=0
\]
由于 \(F\) 在线段 \(AB\) 上，代入直线方程得 \(1+ab=0\)，即 \(ab=-1\). 设 \(R\) 为 \(PQ\) 的中点，则 \(R=\left(-\dfrac12,\dfrac{a+b}{2}\right)\). 于是
\(\overrightarrow{AR}=\left(-\frac{1+a^2}{2},\frac{b-a}{2}\right)\)，\(\overrightarrow{FQ}=(-1,b)\).
由 \(ab=-1\)，有 \(1+a^2=a(a-b)\)，且 \(b=-\dfrac1a\)，所以 \(\overrightarrow{AR}\) 与 \(\overrightarrow{FQ}\) 的对应坐标成比例，故 \(AR\parallel FQ\).

\item 仍设
\(A=\left(\frac{a^2}{2},a\right)\)，\(B=\left(\frac{b^2}{2},b\right)\)，\(P=\left(-\frac12,a\right)\)，\(Q=\left(-\frac12,b\right)\).
记 \(D\) 为直线 \(AB\) 与 \(x\) 轴的交点. 由直线 \(AB\) 的方程，得 \(D=\left(-\dfrac{ab}{2},0\right)\). 在通常的非退化割线约定下，\(l_1\)，\(l_2\) 均不取 \(x\) 轴，故 \(a\)，\(b\ne0\). 以 \(AB\) 为底，利用坐标计算三角形面积，可得
\(S_{\triangle ABF}=\frac12|b-a|\,|FD|\)，\(S_{\triangle PQF}=\frac12|b-a|\).
由题设 \(S_{\triangle PQF}=2S_{\triangle ABF}\)，得 \(|FD|=\dfrac12\). 因为 \(F=\left(\dfrac12,0\right)\)、\(D=\left(-\dfrac{ab}{2},0\right)\)，所以
\[
\left|\frac{1+ab}{2}\right|=\frac12
\]
即 \(ab=0\) 或 \(ab=-2\). 按上述非退化约定舍去 \(ab=0\)，故 \(ab=-2\).

设 \(AB\) 的中点为 \(M(x,y)\)，则
\(x=\frac{a^2+b^2}{4}\)，\(y=\frac{a+b}{2}\).
于是
\[
y^2=\frac{a^2+2ab+b^2}{4}=x+\frac{ab}{2}=x-1
\]
反之，曲线 \(y^2=x-1\) 上任一点都可取实数 \(a\)，\(b\) 使 \(a+b=2y\)、\(ab=-2\)，于是可由上述两点得到该中点. 因此，在非退化约定下，所求轨迹方程为 \(y^2=x-1\).

需要说明的是，原题文字本身没有明说 \(l_1\)，\(l_2\) 不能取 \(x\) 轴. 若严格按字面允许 \(a=0\) 或 \(b=0\)，则 \(ab=0\) 也满足上述面积条件；此时 \(AB\) 的中点满足 \(x=\dfrac{b^2}{4}\)、\(y=\dfrac b2\)（或交换 \(a\)，\(b\)），从而还得到分支 \(y^2=x\). 公开参考解答均按两条直线为非退化割线的通常约定取 \(y^2=x-1\)，本解采用该标准约定并保留这一题面歧义.
\end{enumerate}
\end{solution}

\begin{problem}
设函数 \(f(x)=\alpha\cos 2x+(\alpha-1)(\cos x+1)\)，其中 \(\alpha>0\)，记 \(|f(x)|\) 的最大值为 \(A\).

\begin{enumerate}
\item 求 \(f'(x)\)；
\item 求 \(A\)；
\item 证明 \(|f'(x)|\le2A\).
\end{enumerate}
\end{problem}

\begin{solution}
\begin{enumerate}
\item 由 \((\cos 2x)'=-2\sin 2x\)、\((\cos x)'=-\sin x\)，得
\[
f'(x)=-2\alpha\sin 2x-(\alpha-1)\sin x
\]

\item 令 \(t=\cos x\)，则 \(-1\le t\le1\)，且
\[
f(x)=\alpha(2t^2-1)+(\alpha-1)(t+1)=2\alpha t^2+(\alpha-1)t-1
\]
记右端为 \(g(t)\). 由于 \(\alpha>0\)，\(g(t)\) 是开口向上的二次函数，其顶点横坐标为
\[
t_0=-\frac{\alpha-1}{4\alpha}=\frac{1-\alpha}{4\alpha}
\]
当 \(0<\alpha<\dfrac15\) 时，\(t_0>1\)，所以 \(g(t)\) 在 \([-1,1]\) 上递减. 又
\(g(-1)=\alpha\)，\(g(1)=3\alpha-2\)
且 \(2-3\alpha>\alpha\)，故 \(A=2-3\alpha\).

当 \(\dfrac15\le\alpha\le1\) 时，\(t_0\in[-1,1]\)，且
\[
g(t_0)=-\frac{\alpha^2+6\alpha+1}{8\alpha}
\]
记 \(B=\dfrac{\alpha^2+6\alpha+1}{8\alpha}\). 端点处的绝对值分别为 \(\alpha\) 和 \(|3\alpha-2|\). 直接计算得
\[
B-\alpha=\frac{(1-\alpha)(7\alpha+1)}{8\alpha}\ge0
\]
当 \(\dfrac15\le\alpha\le\dfrac23\) 时，
\[
B-(2-3\alpha)=\frac{(5\alpha-1)^2}{8\alpha}\ge0;
\]
当 \(\dfrac23\le\alpha\le1\) 时，
\[
B-(3\alpha-2)=\frac{(1-\alpha)(23\alpha+1)}{8\alpha}\ge0
\]
因此此时 \(A=B=\dfrac{\alpha^2+6\alpha+1}{8\alpha}\).

当 \(\alpha\ge1\) 时，仍有 \(t_0\in[-1,1]\)，而
\(B-(3\alpha-2)=\frac{(1-\alpha)(23\alpha+1)}{8\alpha}\le0\)，\(3\alpha-2\ge\alpha\).
所以此时 \(A=3\alpha-2\). 综合可得
\[
A=\begin{cases}
2-3\alpha,&0<\alpha<\dfrac15,\\[4pt]
\dfrac{\alpha^2+6\alpha+1}{8\alpha},&\dfrac15\le\alpha\le1,\\[8pt]
3\alpha-2,&\alpha\ge1.
\end{cases}
\]

\item 当 \(0<\alpha<\dfrac15\) 时，由三角不等式
\[
|f'(x)|\le2\alpha|\sin2x|+(1-\alpha)|\sin x|\le1+\alpha
\]
又 \(1+\alpha\le4-6\alpha=2A\)，故 \(|f'(x)|\le2A\).

当 \(\dfrac15\le\alpha\le1\) 时，
\[
|f'(x)|\le2\alpha+(1-\alpha)=1+\alpha
\]
而
\[
2A-(1+\alpha)=\frac{\alpha^2+6\alpha+1}{4\alpha}-(1+\alpha)=\frac{(1-\alpha)(3\alpha+1)}{4\alpha}\ge0
\]
所以 \(|f'(x)|\le2A\).

当 \(\alpha\ge1\) 时，
\[
|f'(x)|\le2\alpha+(\alpha-1)=3\alpha-1\le6\alpha-4=2A
\]
综上，恒有 \(|f'(x)|\le2A\).
\end{enumerate}
\end{solution}

\section{选考题：解答题}

\begin{problem}
如图，\(\odot O\) 中 \(\overset{\frown}{AB}\) 的中点为 \(P\)，弦 \(PC\)，\(PD\) 分别交 \(AB\) 于 \(E\)，\(F\) 两点.

\begin{enumerate}
\item 若 \(\angle PFB=2\angle PCD\)，求 \(\angle PCD\) 的大小；
\item 若 \(EC\) 的垂直平分线与 \(FD\) 的垂直平分线交于点 \(G\)，证明：\(OG\perp CD\).
\end{enumerate}

\begingroup
\FigureLayoutDeclare{img/2016/national_paper_3_science/q22_fig1.png}{0.38\linewidth}{0bp 0bp 0bp 0bp}
\bitmapfigure[float=inline,max width=0.38\linewidth]{img/2016/national_paper_3_science/q22_fig1.png}
\endgroup
\end{problem}

\begin{solution}
\begin{enumerate}
\item 设圆周上各相邻弧 \(AP\)、\(PB\)、\(BD\)、\(DC\)、\(CA\) 所对的圆心角依次为 \(u\)，\(u\)，\(v\)，\(w\)，\(s\)，并记 \(\angle PCD=\theta\). 由圆周角定理，
\[
2\theta=u+v
\]
两弦 \(PD\) 与 \(AB\) 在点 \(F\) 相交，由相交弦所成角的度数等于所对两弧度数和的一半，得
\[
2\angle PFB=u+s+w
\]
由题设 \(\angle PFB=2\angle PCD=2\theta\)，所以
\[
u+s+w=4\theta=2u+2v
\]
即 \(s+w=u+2v\). 又整圆的圆心角和为 \(2u+v+w+s=360^\circ\)，从而
\[
3u+3v=360^\circ
\]
所以 \(u+v=120^\circ\)，故
\[
\angle PCD=\theta=60^\circ
\]

\item 以 \(O\) 为原点，取圆的半径为 \(1\)，并令 \(P=(0,1)\). 由于 \(P\) 是弧 \(AB\) 的中点，弦 \(AB\) 可设为 \(y=k\)，其中 \(-1<k<1\). 令 \(h=1-k\)，则 \(AB\) 上的点可写为 \(E=(e,k)\)、\(F=(f,k)\). 直线 \(PE\) 与单位圆的另一个交点为 \(C\). 设
\[
P+t(E-P)=(te,1-th)
\]
代入单位圆方程得
\[
t^2(e^2+h^2)-2ht=0
\]
除去对应于点 \(P\) 的根 \(t=0\)，得 \(t=\dfrac{2h}{e^2+h^2}\)，所以
\[
C=\left(\frac{2he}{e^2+h^2},\frac{e^2-h^2}{e^2+h^2}\right)
\]
同理，令 \(L_f=f^2+h^2\)，有
\[
D=\left(\frac{2hf}{L_f},\frac{f^2-h^2}{L_f}\right)
\]

设 \(G=(X,Y)\). 记 \(s^2=1-k^2=h(2-h)\). 由上式可得
\[
C-E=\frac{s^2-e^2}{e^2+h^2}(e,-h)
\]
又因为 \(|C|=1\)，所以 \(GE=GC\) 等价于
\[
2G\cdot(C-E)=1-|E|^2=s^2-e^2
\]
在非退化情形下约去公因式，得到
\[
eX-hY=\frac{e^2+h^2}{2}
\]
同理，由 \(GF=GD\) 得
\[
fX-hY=\frac{f^2+h^2}{2}
\]
两式相减，得到
\(X=\frac{e+f}{2}\)，\(Y=\frac{ef-h^2}{2h}\).

令 \(L_e=e^2+h^2\)，则
\[
C-D=\left(-\frac{2h(e-f)(ef-h^2)}{L_eL_f},\frac{2h^2(e-f)(e+f)}{L_eL_f}\right)
\]
因此
\[
\overrightarrow{OG}\cdot\overrightarrow{CD}
=\frac{e+f}{2}\left(-\frac{2h(e-f)(ef-h^2)}{L_eL_f}\right)+\frac{ef-h^2}{2h}\left(\frac{2h^2(e-f)(e+f)}{L_eL_f}\right)=0
\]
所以 \(OG\perp CD\).
\end{enumerate}
\end{solution}

\begin{problem}
在直角坐标系 \(xOy\) 中，曲线 \(C_1\) 的参数方程为
\(\begin{cases}
x=\sqrt3\cos\alpha,\\
y=\sin\alpha
\end{cases}\)
（\(\alpha\) 为参数）；以坐标原点为极点，以 \(x\) 轴的正半轴为极轴，建立极坐标系，曲线 \(C_2\) 的极坐标方程为 \(\rho\sin\left(\theta+\dfrac\pi4\right)=2\sqrt2\).

\begin{enumerate}
\item 写出 \(C_1\) 的普通方程和 \(C_2\) 的直角坐标方程；
\item 设点 \(P\) 在 \(C_1\) 上，点 \(Q\) 在 \(C_2\) 上，求 \(|PQ|\) 的最小值及此时 \(P\) 的直角坐标.
\end{enumerate}
\end{problem}

\begin{solution}
\begin{enumerate}
\item 由参数方程消去参数 \(\alpha\)，得
\[
\frac{x^2}{3}+y^2=\cos^2\alpha+\sin^2\alpha=1
\]
所以 \(C_1\) 的普通方程为 \(\dfrac{x^2}{3}+y^2=1\). 又由极坐标与直角坐标的关系 \(x=\rho\cos\theta\)、\(y=\rho\sin\theta\)，有
\[
\rho\sin\left(\theta+\frac\pi4\right)=\frac{x+y}{\sqrt2}=2\sqrt2
\]
所以 \(C_2\) 的直角坐标方程为 \(x+y=4\).

\item 对任意 \(P=(x,y)\in C_1\)，由柯西不等式
\[
(x+y)^2=\left(\frac{x}{\sqrt3}\cdot\sqrt3+y\right)^2\le\left(\frac{x^2}{3}+y^2\right)(3+1)=4
\]
因此 \(x+y\le2\). 直线 \(C_2\) 为 \(x+y=4\)，所以点 \(P\) 到直线 \(C_2\) 的距离满足
\[
|PQ|\ge\frac{4-(x+y)}{\sqrt2}\ge\frac{2}{\sqrt2}=\sqrt2
\]
当且仅当柯西不等式取等号时，等号条件为
\[
\frac{x}{\sqrt3}:y=\sqrt3:1
\]
结合 \(\dfrac{x^2}{3}+y^2=1\)，得 \(P=\left(\dfrac32,\dfrac12\right)\). 从该点向直线 \(x+y=4\) 作垂线，垂足为 \(Q=\left(\dfrac52,\dfrac32\right)\)，此时 \(|PQ|=\sqrt2\). 故 \(|PQ|\) 的最小值为 \(\sqrt2\)，此时 \(P\) 的坐标为 \(\left(\dfrac32,\dfrac12\right)\).
\end{enumerate}
\end{solution}

\begin{problem}
已知函数 \(f(x)=|2x-a|+a\).

\begin{enumerate}
\item 当 \(a=2\) 时，求不等式 \(f(x)\le6\) 的解集；
\item 设函数 \(g(x)=|2x-1|\)，当 \(x\in\R\) 时，\(f(x)+g(x)\ge3\)，求 \(a\) 的取值范围.
\end{enumerate}
\end{problem}

\begin{solution}
\begin{enumerate}
\item 当 \(a=2\) 时，\(f(x)=|2x-2|+2\). 不等式 \(f(x)\le6\) 等价于
\[
|2x-2|\le4
\]
即 \(-4\le2x-2\le4\). 解得 \(-1\le x\le3\)，所以解集为 \([-1,3]\).

\item 令 \(t=2x\)，则 \(t\in\R\)，题设条件等价于
\[
|t-a|+|t-1|+a\ge3\quad\text{对任意 }t\in\R
\]
两点 \(a\)、\(1\) 到任意点 \(t\) 的距离之和的最小值为 \(|a-1|\)，当 \(t\) 位于 \(a\) 与 \(1\) 之间时取得. 因而题设等价于
\[
a+|a-1|\ge3
\]
当 \(a<1\) 时，左端等于 \(1<3\)，不满足条件；当 \(a\ge1\) 时，左端等于 \(2a-1\)，故 \(2a-1\ge3\)，即 \(a\ge2\). 所以 \(a\) 的取值范围为 \([2,+\infty)\).
\end{enumerate}
\end{solution}
